\subsection{The relative period of the resolved-tail degree-three class}

\begin{proposition}[AT18: Proved; the exact period and its dependence on lifts]
\label{AT18}
Use the resolution, notation, and actual corrected class
$\zeta_{\mathrm{tw}}$ of A-S09-22 over
$R=\mathbb Z_{(2)}$, with $L=I_1\langle s\rangle$.
Thus
\[
 B=aH-Cr-gx,\qquad dB=2\zeta_{\mathrm{tw}},
 \qquad TC=2(K-H),\qquad TH=2J.
 \tag{AT.90}
\]
Let $f:Y^{\mathrm{tf}}\to A$ be one of the actual AT16 target
maps, and put
\[
 \mathcal C_A=\operatorname{Hom}_{sR\text{-}\mathrm{Mod}}(K(2),A).
\]
All products below are products in simplex two; their faces
vanish, so they specify elements of this cochain complex.
Assume the finite-range period comparison used in AT16 holds
in all weights at least $s+1$ through form degree four.
Write
\[
 \kappa_w:H^2(\mathcal C_A^w)\xrightarrow{\;\cong\;}S_R^w
 \tag{AT.91}
\]
for its coefficient-period isomorphism.

The exact obstruction to extending $f$ over the relative two-cell
repair
\[
 du'=\zeta_{\mathrm{tw}},\qquad dv'=2u'-B,
 \qquad |u'|=2,\quad |v'|=1
 \tag{AT.92}
\]
on the sphere parameter $K(2)$ is
\[
 \beta_f=\overline\kappa_{2s+5}
             \bigl([\,f(B)\bmod2\,]\bigr)
       \ \in\ S_R^{2s+5}/2S_R^{2s+5}.
 \tag{AT.93}
\]
Here $\overline\kappa$ is the actual reduction of (AT.91).
The extension exists exactly when $\beta_f=0$.

Suppose additionally that (AT.91) is available in weight $s$.
Then $\beta_f$ is independent of every choice in the AT16
target lift, while the map on the old algebra $X$ is fixed.
In particular, this freedom in the target nullhomotopy does
not change this particular residue.

If the target twist is zero, there is a lift with $\beta_f=0$.
The same conclusion holds whenever actual elements
$U\in\mathcal C_A^{s,3}$ and $V\in\mathcal C_A^{s,2}$ satisfy
\[
               dU=f(z_0),\qquad
               dV=2U-f(\eta_0).
 \tag{AT.94}
\]
Under the additional weight-$s$ hypothesis, every lift then
has zero residue. No nonzero value of (AT.93) in an admissible
target is asserted here.
\end{proposition}

\begin{proof}
The terms of $\mathcal C_A$ are torsionfree: in each form
degree they are subgroups of $A_2$.
The exact reduction sequence and the finite-range comparison
therefore give
\[
 H^2(\mathcal C_A^{w}/2\mathcal C_A^{w})
       \cong H^2(\mathcal C_A^w)/2
       \cong S_R^w/2S_R^w
 \qquad(w\geq s+1),
 \tag{AT.95}
\]
because $H^3(\mathcal C_A^w)=0$.
Since $d f(B)=2f(\zeta_{\mathrm{tw}})$, the element in
(AT.93) is a well-defined cohomology class under this
specific comparison.

Put $w=2s+5$. Vanishing of $H^3(\mathcal C_A^w)$ supplies
$U'\in\mathcal C_A^{w,2}$ with
$dU'=f(\zeta_{\mathrm{tw}})$.
The two-form
\[
                         2U'-f(B)
 \tag{AT.96}
\]
is closed. Its period modulo two is the negative of (AT.93),
which is the same element modulo two. If an actual
$V'\in\mathcal C_A^{w,1}$ has differential (AT.96), this
period is zero.

Conversely, when (AT.93) is zero, the integral period of
(AT.96) is $2c$ for some $c\in S_R^w$.
Choose a closed relative two-form $\Omega$ with period $c$,
using the surjectivity of (AT.91).
Replace $U'$ by $U'-\Omega$. The new expression (AT.96)
has period zero, so the isomorphism (AT.91) says it is a
boundary in $\mathcal C_A$. Choose its one-form primitive
$V'$. These are actual $K(2)$-parameter generator images
satisfying both equations (AT.92); the free relative PD
universal property extends them to a strict map.
This proves both implications of the criterion, with the
displayed morphism.

We next establish the precise invariance needed for the
choice assertion. Consider two actual lift cochains on the
resolved generator complex whose difference is a cochain
homotopy boundary. Work modulo two and denote this homotopy
by $k$, of degree minus one. Put
\[
 \alpha=k(g),\quad \beta=k(x),\quad
 \rho=k(r),\quad \chi=k(a),\quad \nu=k(E),
 \qquad E=u_{e_1}.
\]
The first two have form degree zero, the next two degree
one, and $\nu$ degree two. Let unprimed letters denote the
first target values. The signs in A-S09-22 reduce to
\[
 \begin{aligned}
 \Delta C&=T\beta,& \Delta H&=T\alpha,\\
 \Delta g&=d\alpha+T\rho,&
 \Delta x&=d\beta+T\chi,\\
 \Delta r&=d\rho+T\nu,&
 \Delta a&=d\chi.
 \end{aligned}                                                    \tag{AT.97}
\]
These identities use the differential of the linear generator
complex $W$, not the affine twist: the difference of two
lifts has zero twist. The bottom component of a degree-minus-one
homotopy is zero, and the terminal equation is exactly the
identity for $\Delta C,\Delta H$.

In the target reduced modulo two one has
\[
 dC=Tx,\quad dH=Tg,\quad dg=Tr,\quad dx=Ta,\qquad
 TC=0,\quad T^2=0.
 \tag{AT.98}
\]
The first four equations involve only generator summands on
which the twist is zero. The last two are the actual relations
$TC=2(K-H)$ and $T^2=2T^{[2]}$.
Substitute (AT.97) in $B=aH+Cr+gx$ modulo two and use
(AT.98). Direct expansion gives
\[
 \begin{split}
 B_{\mathrm{new}}-B_{\mathrm{old}}
   =d\bigl(&\chi H+C\rho+g\beta+\alpha x\\
            &+\alpha\,d\beta+T\alpha\chi+T\beta\rho\bigr).
 \end{split}                                                     \tag{AT.99}
\]
For completeness, the cross terms are the differential of
the first four summands plus $CT\nu$; the latter is zero.
The terms quadratic in the homotopy are the differential
of the final three summands, with the remaining terms
multiples of $T^2$. Thus (AT.99) uses no division and proves
equality of the residue for these two lifts.

We now account for every lift choice, rather than only
homotopy boundaries. The finite-range calculation of AT16,
one total degree lower than its obstruction calculation,
identifies the difference classes as
\[
 H^0\operatorname{Hom}(P\otimes K(2),A)
       \cong \operatorname{Ext}^1_{\mathrm{gr}S_R}(L/2L,S_R).
 \tag{AT.100}
\]
All free generators of $P$ have weight at least $s+1$,
so the stated period range applies to every term on the
finite diagonal. The proof is the same actual free-resolution
Hom comparison as in AT16; no unbounded diagonal is used.
The long exact sequence of
$0\to L\xrightarrow{\,2\,}L\to L/2L\to0$ and AT04 give
\[
 \operatorname{Ext}^1_{\mathrm{gr}S_R}(L/2L,S_R)
       \cong S_R^s/2S_R^s.
 \tag{AT.101}
\]
Here $\operatorname{Hom}_{\mathrm{gr}S_R}(L,S_R)=S_R^s$
and all its positive Ext groups vanish.

Use the additional weight-$s$ period comparison to choose
a closed $\omega\in\mathcal C_A^{s,2}$ with
$\kappa_s([\omega])=T^{[s]}$.
There is an actual integral cochain map on $W$ that sends
\[
 v_0(q)\longmapsto \varepsilon(q)\omega
 \tag{AT.102}
\]
and sends all other generator summands to zero.
The equations $\varepsilon\partial=0$ check every
differential, including the bottom quotient. Under (AT.100),
this map represents the generator of (AT.101): its period
is the coefficient multiplication map
$L=I_1\langle s\rangle\to S_R$ by $T^{[s]}$.
Its multiples therefore account for every difference class.

Adding (AT.102) changes only $a$ among the variables in
$B$, by $T\omega$. Thus its change in $B$ is
\[
                         T\omega H=2\omega J,
 \tag{AT.103}
\]
which is already zero modulo two. After subtracting an
appropriate multiple of (AT.102), the difference of any
two lift cochains is an actual Hom boundary by (AT.100).
The resulting homotopy descends through the truncation:
its form-zero component is zero and its form-one equation
factors the bottom quotient, exactly as in AT16.
Equation (AT.99) now proves the claimed independence of
all lift choices.

If the target twist vanishes, assigning zero to every new
generator is an actual lift. It sends $B$ to zero.
For (AT.94), assign instead
\[
 u_0(q)\longmapsto\varepsilon(q)U,\qquad
 v_0(q)\longmapsto\varepsilon(q)V
 \tag{AT.104}
\]
and send every higher-resolution summand and $W^0$ to
zero. The two identities (AT.94) and
$\varepsilon\partial=0$ verify all differential and
bottom relations. This is an actual lift, and it again
sends $C,H,g,x,r$, hence $B$, to zero.
These constructions prove the final assertions.
\end{proof}
